\par\SourceLeaf{126}
Die Verpflichtungen der W. K. zerfallen in drei Hauptrubriken.

I. Verpflichtungen gegen die jetzt vorhandenen Witwen, eventuell deren Kinder,

II. Verpflichtungen gegen diejenigen Professoren, welche jetzt Theilnehmer der W. K. Gesellschaft sind.

III. Verpflichtungen gegen die künftig beitretenden Professoren.

Ad I. Die Verpflichtung gegen jede einzelne Witwe, ohne minorenne Kinder, hat genau den Werth einer Leibrente für dieselbe und lässt sich daher, wenn man ihr jetziges Alter kennt, (und für Mortalitätstabelle und Zinsfuss eine bestimmte Wahl trifft) genau berechnen. Dass in jedem einzelnen Fall ein Entrepreneur, der für diesen Preis die Verpflichtung auf sich nähme, eine Art Glücksspiel spielt, versteht sich von selbst (und werde ich daher im Folgenden, wo immer wieder dieselbe Erinnerung gemacht werden müsste, dies unterlassen, da dies jedem von einiger mathematischer Bildung bekannt ist), aber der Entrepreneur, der mit einer \textit{sehr grossen Zahl} von Personen denselben Contract schlösse, würde, wenn richtig gerechnet ist, mit moralischer Gewissheit nur ein in Proportion zum Ganzen unerhebliches Schwanken zu erwarten haben.

Wo Kinder vorhanden sind, erleidet die Rechnung eine Modification, behält aber dieselbe Gültigkeit wie im vorigen Falle. Natürlich muss aber das Alter der Kinder auch bekannt sein.

Endlich würde streng genommen bei unsrer Witwenkasse, welche sich wiederverheirathende Witwen ausschliesst, noch eine Modification nöthig sein, die scheinbar\footnote{Es gehört nicht hierher, zu rechtfertigen, warum ich diese Einrichtung nur für \textit{scheinbar} vortheilhaft halte, ich bin aber gern bereit, jedem der sich dafür interessirt, die Gründe anzuzeigen.} zum Vortheil der Witwenkasse ist, aber sich natürlich nicht im Voraus berechnen lässt, jedenfalls praktisch \(=0\) zu setzen ist.

Es lässt sich demnach auch der Gesammtbetrag von I. zu Gelde anschlagen, und einen dem Betrage gleichen Theil des Capital-Vermögens der Kasse muss man als dadurch absorbirt betrachten.

Ad II. Auf ähnliche Weise würde sich auch die Verpflichtung gegen jeden einzelnen Professor, der jetzt verheirathetes Mitglied der Gesellschaft ist, nach Gelde anschlagen, wenn es in unserer Gesellschaft ganz ebenso wäre, wie in denjenigen freiwilligen Gesellschaften, wo der jährliche Beitrag oder das Eintrittsgeld nach dem Alter des eintretenden Ehepaars regulirt wird. Das Unterscheidende einer solchen Gesellschaft von der Unsrigen besteht aber in folgendem.

A. In jener erlischt der Contract, wenn die Frau vor dem Manne stirbt; soll er bei einer Wiederverheirathung erneuert werden, so ist es ein ganz neuer nach dem Alter der zweiten Frau zu regulirender Contract. Bei uns sind auch unverheirathete Mitglieder, die eine Braut in beliebigem Alter wählen können, ebenso Witwer, die möglicher Weise sich wieder verheirathen können. Dies alles kann aber jetzt einer Vorausberechnung gar nicht unterworfen werden. Ich würde aber glauben, dass wenn man diejenigen Mitglieder, die jetzt verheirathet sind, nach ihrem und nach dem Alter ihrer jetzigen Frauen einem Calcül unterzöge, und dann für die übrigen jetzt nicht verheiratheten den Mittelwerth jener ersteren Resultate zum Grunde legte, es sich so ziemlich compensiren würde. Möglicherweise werden von einigen jetzt verheiratheten Mitgliedern nicht ihre jetzigen Frauen sondern zweite oder dritte dereinst die Witwenpension geniessen, dagegen wird aber ohne Zweifel ein Theil der jetzt nicht verheiratheten in diesem Stande bleiben. Ich sehe wenigstens nicht ab, was man hier mehr thun könne, als die zweierlei Eventualitäten, die einen zum Nachtheil, die andern zum Vortheil der Kasse gereichend, sich gegenseitig aufheben zu lassen.

B. Ausserdem findet aber auch noch der Unterschied statt, dass Kinder, vielleicht jetzt noch gar nicht geboren, demnächst möglicherweise, an den Vortheilen Theil nehmen. Auch das lässt sich daher so nicht veranschlagen; ich glaube jedoch, dass für diese Unvollkommenheit sich ein völlig ausreichendes Surrogat finden lässt, welches ich aber um nicht gar zu weitläufig zu werden, hier nicht näher entwickeln will.
